%% FONCTION MATHEMATIQUES DE BASE %%
%
\section{Fonction logistique dans un domaine de d�finition}

\itemstyle{Op�rateur~:} \operfmt{Logistic\_domain}
\index{Logistic\_domain}

%%%%%OBJET%%%%
\itemstyle{Objet~:}
Application d'une fonction logistique (de degr� 2) appliqu�e dans un domaine de d�finition. Si la variable en entr�e n'appartient pas au domaine de d�finition (valeur comprise entre borne mini et borne maxi), alors l'op�rateur renvoie une valeur manquante.

%%%%%ARGUMENTS%%%%
\bigskip
%% Dans ce fichier, les �ventuels accents sont cod�s
%% en ISO-8859-15 (aka latin9)
%% La ligne suivante donne trois exemples 
%% e accent aigu : �, e accent grave : �, e-dans-l'o : �
%%%%%ENTREES%%%%
\begin{tabarg}
\multicolumn{3}{l}{\titretab {Entr\'ees}}\\
\hline
$X$ & Vecteur, {\tt float} & Op\'erande (unit\'e quelconque)\\
$A$ & Coeff., \texttt{double} & Coefficient multiplicateur A du terme exponentiel()\\
$K$ & Coeff., \texttt{double} & Coefficient multiplicateur K de la  fonction()\\
$\alpha$ & Coeff., {\tt double} & Coefficient pour l'exposant du terme de degr\'e 0.\\
$\beta$ & Coeff., {\tt double} & Coefficient pour l'exposant du terme de degr\'e 1.\\
$\gamma$ & Coeff., {\tt double} & Coefficient pour l'exposant du terme de degr\'e 2.\\
$Borne\_{min}$ & Coeff., {\tt float} & Borne minimale pour le parametre X\\
$Borne\_{max}$ & Coeff., {\tt float} & Borne maximale pour le parametre X\\
$Unit\_{res}$ & Coeff., {\tt String} & Unit\'e du r\'esultat.\\
\end{tabarg}

\medskip
%%%%%SORTIES%%%%
\begin{tabarg}
\multicolumn{3}{l}{\titretab {Sorties}}\\
\hline
$Z$ & Vecteur, \texttt{float} & R\'esultat (unit\'e indiqu\'ee par $Unit\_{res}$) \\
\end{tabarg}


%%%%%FORMULE%%%%
\itemstyle{Formule}

%% Formule de l'operateur : Logistic_domain
%% Dans ce fichier, les �ventuels accents sont cod�s
%% en ISO-8859-15 (aka latin9)
%% La ligne suivante donne trois exemples 
%% e accent aigu : �, e accent grave : �, e-dans-l'o : �
\[
Z = \frac{K}{1+A \times e^ {\alpha + \beta X +\gamma X^{2i}}}
\]



%%%%%COMMENTAIRE%%%%
\itemstyle{Commentaire~:}
On effectue le calcul uniquement si la variable X est dans le domaine de d�finition :
 c'est � dire sup�rieure ou �gale � Borne\_mini et inf�rieure ou �gale � Borne\_maxi.


%%%%%AUTEUR%%%%
\itemstyle{Auteur~:}
CNRM/GMEI/TRAMM

%%%%%REFERENCE%%%%
\itemstyle{R�f�rences~:}
